\chapter{Увод}
\label{ch:uvod}

Съвременната музикална продукция е немислима без програмиране на барабанни
партии в цифрова аудио работна станция (\eng{Digital Audio Workstation, DAW}).
Когато обаче партията бъде наредена „на ръка'' (с мишка) всяка нота
и с еднаква сила на удара, резултатът звучи механично, което често може да бъде нежелано.
Липсва ѝ онова, което слушателят възприема като
„човешко'' изпълнение: вариацията на момента на изсвирване, както и -
не по-малко важно - на силата на изсвирване (\emph{динамиката} на изпълнението). 
В различни софтуери има вече внедрени решения за първия проблем (т.нар \eng{humanize}),
но изследване и разработки относно втория проблем почти липсват.
Настоящата работа се занимава именно с това —
автоматичното \emph{предсказване на динамиката} на
барабанна партия в MIDI формат.

\section{Преглед на проблемната област}
\label{sec:pregled}

В стандарта MIDI, силата, с която е ударена дадена нота, се кодира чрез
целочислена стойност \eng{velocity} в диапазона $0$--$127$ включително. За барабаните тази
стойност не е просто въпрос на „по-силно'' или „по-тихо'': тя определя
акцентите и т.нар.\ „призрачни ноти'' (\eng{ghost notes}), оформя усещането за
ритъм (\eng{groove}) и носи голяма част от стилистичната идентичност на
изпълнителя. Един и същ ритмичен рисунък, изсвирен от човек барабанист, показва
непрекъснато вариране на \eng{velocity} между ударите; същият рисунък, въведен
програмно с постоянна сила, звучи изкуствено и безжизнено.

Съществуващите подходи за придаване на човешко звучене се разделят условно на
няколко групи: системи, базирани на правила, които прилагат детерминистични
евристики върху времето и силата~\cite{friberg2006overview}; статистически
модели, които извличат закономерности от записи на реални изпълнения; и
съвременни невронни модели за генериране на груув, като
\eng{GrooVAE}~\cite{gillick2019groove}. Ключовото наблюдение, което мотивира
тази работа, е, че почти цялото внимание в литературата и в практическите
инструменти е насочено към придаване на човешко звучене на \emph{времето}
(микротайминг, суинг), докато \emph{динамиката} остава слабо изследвана и разработена като
самостоятелна задача.

\section{Мотивация}
\label{sec:motivacia}

Задачата е практически значима. Музикант или композитор, който не разполага 
с достатъчно чувствителни електронни барабани или MIDI клавиатура, след като нареди 
барабанна партия на мрежа, се нуждае от способ, който да ѝ придаде правдоподобна динамика,
съобразена с желания от него/нея стил.

Това повдига и интересен научен въпрос: доколко изобщо човешката динамика е
\emph{предсказуема} само от структурна и времева информация, при пълна забрана да се
използват съседните на въпросната нота стойности за \eng{velocity} (силата на удара)? Ако значителна част от вариацията
е обяснима, това е пряко приложимо в DAW-и и драм машини; ако не е, самата
степен на непредсказуемост е информативна за това какво прави изпълнението
човешко.

\section{Цели на дипломната работа}
\label{sec:celi}

Целите на настоящата дипломна работа са:
\begin{itemize}
  \item Да се формализира задачата за предсказване на човешката динамика
        (\eng{velocity}) на ноти на ударни инструменти въз основа единствено на структурни и
        времеви признаци.
  \item Да се проектират и обучат различни подходящи модели, които
        предсказват динамиката на дадени ноти при подходящи входни данни.
  \item Да се проведе цялостна експериментална оценка върху корпус от данни
        E-GMD с метрики, надхвърлящи простата средна абсолютна грешка, включващи
        мерки за запазване на динамичния обхват и за коректност на акцентите.
  \item Да се изследва устойчивостта на модела към нечуван изпълнител, както и
        да се направи анализ на интерпретируемостта и грешките, свързвайки
        наученото с реални явления при изпълнението на барабани.
  \item Да се реализира софтуерна приставка за популярна цифрова аудио работна станция
\end{itemize}

\section{Структура на дипломната работа}
\label{sec:struktura}

Глава~\ref{ch:osnova} представя необходимата теоретична основа — възприятието
на барабанната динамика, представянето на партии в MIDI, признаците от
структура и време, както и корпуса от данни E-GMD. Глава~\ref{ch:svarzani}
разглежда свързаните разработки и позиционира приноса на работата спрямо тях.
Глава~\ref{ch:metod} формализира задачата и описва метода — извличането на
признаци и семейството от модели. Глава~\ref{ch:eksperimenti} представя
експерименталната постановка, метриките и резултатите, включително анализ на
устойчивост и интерпретируемост. Глава~\ref{ch:realizacia} описва програмната
реализация. Глава~\ref{ch:zaklyuchenie} обобщава постигнатото и очертава насоки
за бъдещо развитие.
